\section{Research Question}

\begin{itshape}
Does replacing joint household income taxation with individual income taxation reduce the gender gap in labor supply and human capital accumulation among dual-earner couples?
\end{itshape}
\noindent Under joint taxation the secondary earner's income is taxed at the household's top marginal rate, creating a higher effective marginal tax rate on female labor than on male labor even when the statutory schedule is gender-neutral.
I ask whether this tax wedge is a quantitatively important driver of the gender gap in hours worked and, through the human capital channel, of the gender wage gap over the life cycle.
The question is directly policy-relevant: Germany and France retain joint taxation or income-splitting systems, while Denmark abolished joint taxation in 1983 --- offering a natural experiment to discipline the model \citep{bick2017taxation}.\vspace{-1em}
